% ═══════════════════════════════════════════════════════════════════════
%  Stratégie 2 --- XAUUSD niveaux x10. Squelette : aucun chiffre tant que
%  la campagne n'a pas tourné. Les résultats arrivent par les macros de
%  tables/x10_headline.tex.
% ═══════════════════════════════════════════════════════════════════════
\section{Glossaire}
\label{app:x10_glossaire}

% TODO L2 : compléter par les termes que la prose finale introduira
% (plateau, essai, gel), en gardant une définition par terme, sans chiffre.

Les définitions ci-dessous sont celles retenues \emph{dans ce document}. Deux
d'entre elles --- le \code{VWAP} et le \code{DXY} --- s'écartent de l'usage le
plus répandu~; l'écart est signalé et motivé.

\subsection{Instrument et unités de temps}

\begin{description}
  \item[Niveau x10] Prix multiple de dix dollars sur \code{XAUUSD}. Ces
  niveaux forment une grille régulière, fixée une fois pour toutes et
  indépendante du marché~: c'est le seul repère sur lequel la stratégie
  travaille.
  \item[M5] Barre de cinq minutes. C'est l'unité de décision~: signaux,
  entrées et sorties y sont évalués.
  \item[H1] Barre d'une heure. C'est l'unité de contexte~: elle oriente, elle
  ne déclenche pas. Toute série H1 est décalée d'une barre avant d'être lue en
  M5, pour qu'aucune information ne soit utilisée avant d'exister.
\end{description}

\subsection{Contexte}

\begin{description}
  \item[EMA50] Moyenne mobile exponentielle sur cinquante barres H1. Donne la
  direction dominante du contexte.
  \item[VWAP] \emph{Ici}~: moyenne cumulée de séance du prix typique
  $(H+L+C)/3$, \emph{non pondérée par les volumes}. L'usage courant pondère par
  le volume échangé~; l'or au comptant n'en publie pas, et les volumes
  rapportés par les courtiers sont des décomptes de ticks propres à chacun.
  Une moyenne non pondérée est identique sur les trois moteurs, ce qu'une
  moyenne pondérée ne serait pas.
  \item[DXY] \emph{Ici}~: panier synthétique à quatre jambes ---
  \code{EURUSD}, \code{USDJPY}, \code{GBPUSD}, \code{USDCAD} --- aux poids ICE
  renormalisés. L'indice officiel en compte six~; les deux jambes écartées sont
  indisponibles ou illiquides sur les trois plateformes. Le panier mesure la
  force du dollar, qui est ce dont la stratégie a besoin.
  \item[ATR] \emph{Average True Range}~: amplitude moyenne des barres récentes.
  Sert d'unité de mesure locale~: toute distance ou vitesse est divisée par
  l'ATR pour rester comparable d'une année à l'autre.
\end{description}

\subsection{Approche du niveau}

\begin{description}
  \item[Vitesse] Distance parcourue par le prix par unité de temps à
  l'approche du niveau, exprimée en ATR.
  \item[Accélération] Variation de cette vitesse. Une approche qui accélère et
  une approche qui s'essouffle ne présagent pas du même dénouement.
  \item[Breakout] Franchissement du niveau dans le sens de l'approche, suivi
  d'une poursuite du mouvement.
  \item[Sweep] Franchissement bref, qui va chercher les ordres placés
  au-delà du niveau puis revient. Distinguer un sweep d'un breakout est le
  problème central de la stratégie.
  \item[Réintégration] Retour du prix du côté d'où il venait après un
  franchissement. C'est la confirmation d'un sweep et la condition d'entrée des
  scénarios de retournement.
\end{description}

\subsection{Mesure et validation}

\begin{description}
  \item[R] Unité de risque~: la perte qu'un trade subit si son stop est
  touché. Un gain de deux~R vaut deux fois le risque engagé. Le mandat exige
  une géométrie d'au moins $R \geq 1$.
  \item[In-sample] Période sur laquelle la stratégie a été conçue et ses
  paramètres choisis. Les résultats y sont, par construction, optimistes.
  \item[Hors échantillon] Période tenue à l'écart de toute décision de
  conception, lue une seule fois, après gel. C'est la seule mesure qui ait
  valeur de test.
  \item[DSR] \emph{Deflated Sharpe Ratio}~: Sharpe corrigé du nombre d'essais
  menés pour l'obtenir. Chercher longtemps produit mécaniquement de beaux
  Sharpes~; le DSR retire cette part de chance.
  \item[PBO] \emph{Probability of Backtest Overfitting}~: probabilité que la
  configuration retenue en période d'apprentissage se classe dans la moitié
  basse en période de contrôle. Une valeur élevée signale une sélection qui ne
  survit pas au changement d'échantillon.
\end{description}
